% Appendix D --- Simplicity on Elements: the dominant-corner existence proof.
% Content lifted from earlier §4.5 "Simplicity on Elements: A
% Cross-Substrate Reference" when §4 was compressed to main-body
% measurements only. Simplicity is not a Bitcoin covenant
% proposal and is excluded from the formal impossibility results in §3;
% it appears here as a constructive existence proof that the dominant
% corner D^\star is occupiable on some substrate.

\section{Simplicity on Elements: The Unoccupied Optimal Corner}
\label{app:simplicity}

\subsection{Scope and motivation}
\label{app:simplicity:scope}

Simplicity~\cite{OC17} is an \emph{alternative script primitive}---a
typed combinator calculus with Coq-verified denotational
semantics---deployed on Elements~\cite{DPNS16}, a Bitcoin-adjacent
sidechain whose blocks are produced by known functionaries rather
than by proof-of-work miners. Execution uses native-code
\emph{jets} for known combinator subtrees. Because Elements'
consensus model differs from Bitcoin's, we exclude Simplicity from
the formal results of Section~\ref{sec:impossibility} and report it
here as a cross-substrate reference: Simplicity demonstrates that
the dominant corner $D^\star$ of
Section~\ref{sec:imposs:design_space} is occupiable on some
substrate, even while no Bitcoin covenant proposal
occupies it.

This distinction is load-bearing. The claim of the paper is
\emph{unoccupied on Bitcoin}, not \emph{theoretically
unoccupiable}. Simplicity is the constructive existence proof for
the latter.

\subsection{Design-space position}
\label{app:simplicity:axes}

Our Simplicity vault occupies the strict-best axis values on every
axis where such a value exists:

\begin{itemize}
  \item \axauth{} = \texttt{cold-key}: key-gated recovery via a
    Schnorr signature check.
  \item \axrevault{} = \texttt{no}: atomic withdrawal; the program
    enforces the complete output tuple via
    \texttt{jet::outputs\_hash()\,=\,param::VAULT\_RECOVER\_OUTPUTS\_HASH}
    (see \texttt{simf/vault.simf}).
  \item \axfee{} = \texttt{dynamic} (Elements fee structure):
    fees live inside the transaction; no out-of-band anchor.
  \item \axwdbind{} = \texttt{dedicated-jets-combinator}: the
    withdrawal destination is bound through a typed output-hash
    commitment at compile time, stronger than any of the
    Bitcoin proposals on the same axis.
  \item \axrecbind{} = \texttt{dedicated-jets-combinator}: same
    binding for the recovery leaf, eliminating
    class \classcold{}.
  \item \axopsurf{} = \texttt{single-mode}: a program's meaning is
    determined by its Merkle root; no \texttt{OP\_SUCCESSx}-style
    fallthrough surface.
  \item \axintro{} = \texttt{typed-combinator-introspection}: the
    fundamental operators are total, typed, and Coq-verified.
\end{itemize}

Consequently, Simplicity's vault program is immune by construction
to classes \classfee{}, \classgrief{}, \classscale{}, \classhot{},
\classcold{}, and \classmode{} --- every attack class identified in
Section~\ref{sec:impossibility}. It occupies the strict-best
position on \axauth{}, \axrevault{}, \axwdbind{}, \axrecbind{}, and
\axintro{} simultaneously; no Bitcoin proposal does.

\subsection{Lifecycle measurements}
\label{app:simplicity:empirical}

The Simplicity vault lifecycle totals 865\vB{} (deposit 269,
trigger 298, withdraw 298, recover 286), $49$--$135\%$ larger than
the Bitcoin designs of Table~\ref{tab:vsizes}. The overhead bundles
Elements' serialisation (${\sim}30$--$40$ bytes per output for
confidential-transaction fields even in non-confidential mode) and
Simplicity's bit-oriented DAG witness format. Elements functionaries
do not operate a public mempool, so class \classfee{} (fee-channel
pinning) has no substrate analogue under honest-federation
assumptions.

Our Simplicity vault is atomic: recovery binds the complete output
tuple via the jets commitment cited above, admitting neither partial
withdrawal nor batched recovery and placing the program at
\axrevault{} = \texttt{no} in the design-space lattice. Class
\classscale{} (per-UTXO recovery-cost scaling) has no surface. Cold-key
compromise (TM11) forces recovery but cannot redirect funds, matching
\opctv{} and \opvault{}. A different Simplicity program could express
\axrevault{} = \texttt{yes} by replacing the whole-tuple commitment
with per-output checks; that design would inherit \classscale{}
under D2 (Section~\ref{sec:imposs:class_scale}) but is not what we
measured.

\subsection{Foundations and substrate}
\label{app:simplicity:foundations}

Simplicity's Coq-verified denotational semantics give stronger
guarantees than Alloy bounded model checking of the Bitcoin Script
extensions can produce --- a Simplicity program's meaning is
determined by its Merkle root, under a formal semantics that has
been mechanically verified. However, Simplicity on Liquid is live
today under \emph{federation trust}, whereas the four Bitcoin Script
extensions remain under \emph{Bitcoin consensus review}. The
substrate choice is a security-relevant decision alongside the
per-covenant analysis of Section~\ref{sec:impossibility}: a vault
on Simplicity inherits the federation's liveness and censorship
properties; a vault on Bitcoin inherits proof-of-work's. The
dominant corner being occupiable on Simplicity is therefore not a
turn-key prescription to migrate vaults to Liquid; it is evidence
that the unoccupied status among current Bitcoin proposals reflects a
consensus-substrate constraint rather than a design impossibility.
